\documentclass[tikz,border=10pt]{standalone}
\usepackage[no-math]{fontspec}
\usepackage[brazil]{babel}
\usepackage{amsmath,amssymb}
\IfFileExists{C:/Windows/Fonts/Lato_Regular.ttf}{%
  \setmainfont{Lato}[
    Path = C:/Windows/Fonts/,
    UprightFont = Lato_Regular.ttf,
    BoldFont = Lato_Bold.ttf,
    ItalicFont = Lato_Italic.ttf,
    BoldItalicFont = Lato_BoldItalic.ttf
  ]
  \setsansfont{Lato}[
    Path = C:/Windows/Fonts/,
    UprightFont = Lato_Regular.ttf,
    BoldFont = Lato_Bold.ttf,
    ItalicFont = Lato_Italic.ttf,
    BoldItalicFont = Lato_BoldItalic.ttf
  ]
}{%
  \RequirePackage[default]{lato}
}
\renewcommand{\familydefault}{\sfdefault}

\usepackage{xcolor}
\usepackage{tikz}
\usepackage{pgfplots}
\pgfplotsset{compat=1.18}
\usetikzlibrary{arrows.meta,positioning,shapes.geometric,fit,backgrounds,calc,decorations.pathreplacing,patterns}

% Paleta institucional UEFS / Bioengenharia de Solos
\definecolor{bsNavy}{HTML}{1F4268}
\definecolor{bsNavyDark}{HTML}{152E4A}
\definecolor{bsGold}{HTML}{B38A3D}
\definecolor{bsGoldDark}{HTML}{8C6927}
\definecolor{bsTeal}{HTML}{0E7C7B}
\definecolor{bsForest}{HTML}{2D6A4F}
\definecolor{bsAlert}{HTML}{9E2A2B}
\definecolor{bsEarth}{HTML}{8C5A3C}
\definecolor{bsWater}{HTML}{2E86AB}
\definecolor{bsInk}{HTML}{26313D}
\definecolor{bsMuted}{HTML}{5B6573}

\definecolor{bgNavy}{HTML}{EDF2F7}
\definecolor{bgGold}{HTML}{F8F3E9}
\definecolor{bgTeal}{HTML}{EBF6F5}
\definecolor{bgForest}{HTML}{EEF5F0}
\definecolor{bgAlert}{HTML}{FDF2F2}

\tikzset{
  >=Stealth,
  bsBox/.style={
    draw=#1, fill=white, rounded corners=5pt, line width=1.15pt,
    inner sep=7pt, align=center, font=\footnotesize\sffamily, text=bsInk
  },
  bsTint/.style={
    draw=#1, fill=#1!8!white, rounded corners=5pt, line width=1.15pt,
    inner sep=7pt, align=center, font=\footnotesize\sffamily, text=bsInk
  },
  bsBadge/.style={
    fill=#1, text=white, rounded corners=9pt, inner xsep=8pt, inner ysep=3pt,
    font=\scriptsize\bfseries\sffamily
  },
  bsTitle/.style={font=\normalsize\bfseries\sffamily, text=bsNavy, align=center},
  bsSub/.style={font=\footnotesize\sffamily, text=bsMuted, align=center},
  bsArrow/.style={->, >=Stealth, line width=1.35pt},
  bsDArrow/.style={<->, >=Stealth, line width=1.35pt}
}


\begin{document}
\begin{tikzpicture}[font=\sffamily]

% ============================================================
% TÍTULO GERAL
% ============================================================
\node[bsTitle, anchor=west] at (0, 5.45)
  {Sequência Construtiva da Parede Krainer em Campo};
\node[bsSub, anchor=west] at (0, 5.00)
  {Etapas executivas da fundação em trincheira até a consolidação vegetal da crista};

% ============================================================
% 5 CARDS DE FASE HORIZONTAIS
% ============================================================
\node[bsTint=bsAlert, text width=3.10cm, minimum height=3.65cm, anchor=north west] (p1) at (0, 4.45) {};
\node[bsBadge=bsAlert, rounded corners=5pt, anchor=north] (b1) at ($(p1.north)+(0,-0.22)$) {FASE 1};
\node[font=\footnotesize\bfseries\sffamily, text=bsNavyDark, align=center, text width=2.95cm, anchor=north] (t1)
  at ($(b1.south)+(0,-0.16)$) {Escavação da Base\\Abaixo do\\Solapamento};
\node[font=\scriptsize\sffamily, text=bsInk, align=center, text width=2.90cm, anchor=north]
  at ($(t1.south)+(0,-0.16)$) {Abrir trincheira abaixo\\da cota de solapamento,\\nivelando o apoio.};

\node[bsTint=bsEarth, text width=3.10cm, minimum height=3.65cm, anchor=north west] (p2) at (3.85, 4.45) {};
\node[bsBadge=bsEarth, rounded corners=5pt, anchor=north] (b2) at ($(p2.north)+(0,-0.22)$) {FASE 2};
\node[font=\footnotesize\bfseries\sffamily, text=bsNavyDark, align=center, text width=2.95cm, anchor=north] (t2)
  at ($(b2.south)+(0,-0.16)$) {Primeira Camada de\\Toras \& Conectores};
\node[font=\scriptsize\sffamily, text=bsInk, align=center, text width=2.90cm, anchor=north]
  at ($(t2.south)+(0,-0.16)$) {Posicionar longarinas e\\travessas com pinos ou\\parafusos na grelha inicial.};

\node[bsTint=bsGold, text width=3.10cm, minimum height=3.65cm, anchor=north west] (p3) at (7.70, 4.45) {};
\node[bsBadge=bsGold, rounded corners=5pt, anchor=north] (b3) at ($(p3.north)+(0,-0.22)$) {FASE 3};
\node[font=\footnotesize\bfseries\sffamily, text=bsNavyDark, align=center, text width=2.95cm, anchor=north] (t3)
  at ($(b3.south)+(0,-0.16)$) {Camada de Solo\\Compactado + Ramos Vivos};
\node[font=\scriptsize\sffamily, text=bsInk, align=center, text width=2.90cm, anchor=north]
  at ($(t3.south)+(0,-0.16)$) {Preencher e compactar\\solo fértil e dispor ramos\\vivos até o solo natural.};

\node[bsTint=bsNavy, text width=3.10cm, minimum height=3.65cm, anchor=north west] (p4) at (11.55, 4.45) {};
\node[bsBadge=bsNavy, rounded corners=5pt, anchor=north] (b4) at ($(p4.north)+(0,-0.22)$) {FASE 4};
\node[font=\footnotesize\bfseries\sffamily, text=bsNavyDark, align=center, text width=2.95cm, anchor=north] (t4)
  at ($(b4.south)+(0,-0.16)$) {Elevação Escalada\\(Inclinação para Dentro)};
\node[font=\scriptsize\sffamily, text=bsInk, align=center, text width=2.90cm, anchor=north]
  at ($(t4.south)+(0,-0.16)$) {Repetir camadas mantendo\\inclinação de 10–15\%\\contra o talude.};

\node[bsTint=bsForest, text width=3.10cm, minimum height=3.65cm, anchor=north west] (p5) at (15.40, 4.45) {};
\node[bsBadge=bsForest, rounded corners=5pt, anchor=north] (b5) at ($(p5.north)+(0,-0.22)$) {FASE 5};
\node[font=\footnotesize\bfseries\sffamily, text=bsNavyDark, align=center, text width=2.95cm, anchor=north] (t5)
  at ($(b5.south)+(0,-0.16)$) {Drenagem de Crista\\\& Acabamento\\Vegetal};
\node[font=\scriptsize\sffamily, text=bsInk, align=center, text width=2.90cm, anchor=north]
  at ($(t5.south)+(0,-0.16)$) {Executar canaleta de\\crista, recobrir o topo e\\plantar mudas nativas.};

% Setas entre fases
\draw[bsArrow, bsNavy] (p1.east) -- (p2.west);
\draw[bsArrow, bsNavy] (p2.east) -- (p3.west);
\draw[bsArrow, bsNavy] (p3.east) -- (p4.west);
\draw[bsArrow, bsNavy] (p4.east) -- (p5.west);

% Nota de controle de qualidade
\node[draw=bsForest!35, fill=bgForest, rounded corners=5pt, line width=0.9pt,
      text width=18.50cm, inner xsep=10pt, inner ysep=6pt, anchor=north west,
      font=\scriptsize\sffamily, text=bsInk] at (0, 0.45)
  {\textbf{\color{bsForest}Boas práticas de campo:} Manter a extremidade basal dos ramos vivos sempre em contato com o solo úmido do maciço e proteger o topo contra enxurradas concentradas durante o primeiro ciclo chuvoso.};

\end{tikzpicture}
\end{document}
